\section{问题提出}

\subsection{问题背景}

神经网络模型的算力需求增长很快，单个 AI 核心受成本、功耗和散热限制，多核 NPU 成为提升吞吐的主要途径。多核执行要把一张完整的计算图切成若干子图，分给不同核心，并确定每个核心上子图的先后顺序。切得越细，可并行的工作越多，但跨子图的数据要经过共享主存 DDR 中转，所有核心争用同一份 60 bytes/cycle 的带宽；切得越粗，数据可以留在核内缓存里复用，但并行度下降，大子图还会因为缓存放不下而产生换入换出。本题要求在这些相互冲突的因素之间找到平衡，使总执行时间（Makespan）尽可能短，同时控制额外的数据搬运。

题目给出的计算图是分块计算展开后的数据流图：张量节点表示一块数据，操作节点表示一次计算或搬运，边只连接张量与操作。每个核心与达芬奇架构\cite{liao2019davinci}类似，有 Cube（PIPE\_M）与 Vector（PIPE\_V）两个计算单元、两条搬运流水线（PIPE\_MTE2、PIPE\_MTE3），以及私有的 L1（524{,}288 字节）和 UB（131{,}072 字节）缓存。参赛算法只需输出两样东西：每个非搬运操作所属的子图编号，以及每个核心上的子图执行序列。之后由官方评估程序自动插入边界搬运、调用给定的核内启发式调度器确定每个 Task 内的操作顺序，并以离散事件方式模拟所有核心在共享带宽下的并行执行，给出 Makespan、总额外搬运量和缓存命中率。

题目区分两种硬件场景。场景 A 没有核间同步机制，每个子图是一个独立 Task，Task 之间不保留核内数据，所有跨子图数据都要写回 DDR 再读入；同核前一个 Task 结束后至少等待 100 周期、每个跨核前驱结束后至少等待 1000 周期才能激活。场景 B 有核间同步机制，同一核心上的全部子图合并成一个 Task，同核子图之间的数据可以驻留在 L1/UB 中直接复用，只有跨核依赖需要经 DDR 中转，并额外付出 500 周期的同步延迟。在场景 B 的基础上，问题三再引入所有核心共享的只读 L2 Cache（容量 1~MB、带宽 250 bytes/cycle），它按先进先出方式缓存 COPY\_IN 读过的张量，带宽与 DDR 相互独立。

\subsection{问题提出}

问题一针对场景 A。要求建立多核切图与调度模型并设计求解算法，为每个非 COPY\_IN/COPY\_OUT 操作分配子图编号，并确定各核心上的子图执行顺序，以缩短 Makespan 为主要目标、兼顾总额外数据搬运量。算法要在 2--5 核范围内都能稳定、高效地给出高质量结果，题目对“高效”的说明是避免暴力迭代搜索，单个用例以 5--10 分钟内给出结果为佳。结果包括三部分：正文给出 1--5 核平均加速比折线图，其中单核基准为核内调度算法在整图上得到的 Makespan，单核加速比为 1，$k$ 核加速比为单核基准与 $k$ 核 Makespan 之比，平均加速比取所有用例加速比的算术平均；附录逐用例列出评估程序给出的 Makespan 与总额外搬运量；附件提供可复现的 Python 程序，输入输出格式遵循题目附录 B。

问题二针对场景 B。要求在问题一的基础上建立适用于场景 B 的模型与算法，在严格满足 L1/UB 容量约束的前提下进一步缩短 Makespan、减少额外搬运，同样要在 2--5 核范围内稳定、高效，提交的正文图、附录表和附件程序与问题一形式相同。问题三在场景 B 的基础上引入共享只读 L2 Cache，要求研究其容量和带宽对切图与调度的影响，并建立模型、设计算法。正文需给出 1--5 核下无 L2 与只读 Cache 两种配置的对比曲线，并报告相同核数下只读 Cache 相对无 L2 的加速比，即无 L2 的 Makespan 与只读 Cache 的 Makespan 之比；附录逐用例列出两种配置的 Makespan、总额外搬运量和 Cache 命中率。全部结果都必须由官方评估程序在固定配置下给出。
